\documentclass[tikz,border=6pt]{standalone}
\usepackage{ctex}
\usepackage{amsmath}
\usepackage{pgfplots}
\pgfplotsset{compat=1.18}
\usetikzlibrary{calc,angles,quotes,arrows.meta,positioning,decorations.markings}
\begin{document}
\begin{tikzpicture}
\begin{axis}[
  width=16cm, height=10cm,
  axis lines=middle,
  xmin=-1.3, xmax=4.6, ymin=-2.6, ymax=4.9,
  xtick={0,0.7853982,1.5708,2.356194,3.14159,3.926991},
  xticklabels={$0$,$\frac{\pi}{4}$,$\frac{\pi}{2}$,$\frac{3\pi}{4}$,$\pi$,$\frac{5\pi}{4}$},
  ytick={-2,-1,1,2,3,4},
  xlabel={$x$}, ylabel={$y$},
  xlabel style={anchor=west}, ylabel style={anchor=south},
  tick label style={font=\footnotesize},
  clip=false,
  legend style={at={(0.015,0.985)},anchor=north west,font=\footnotesize,fill=white,fill opacity=0.85,text opacity=1,draw=gray!50},
  legend cell align=left,
]
  % 基准 y = sin x
  \addplot[gray,thick,dashed,smooth,samples=200,domain=-1.2:4.5] {sin(deg(x))};
  \addlegendentry{基准 $y=\sin x$}
  % 目标 y = 3 sin(2x - pi) + 1
  \addplot[blue,very thick,smooth,samples=300,domain=-1.2:4.5] {3*sin(deg(2*x - 3.14159)) + 1};
  \addlegendentry{$y=3\sin(2x-\pi)+1$}

  % 中线 k = 1
  \addplot[red!70!black,thick,dotted,domain=-1.2:4.5] {1};
  \draw (axis cs:4.45,1) node[right,text=red!70!black,font=\footnotesize]{中线 $k=1$};

  % ===== 振幅 A=3 标注（在峰谷之间，避开峰值）=====
  % 第一个峰在 x=3pi/4=2.356，y=4；第一个谷在 x=pi/4=0.785，y=-2
  \draw[<->,blue!60!black,thick] (axis cs:2.356194,1) -- (axis cs:2.356194,4);
  \draw (axis cs:2.356194,2.55) node[fill=white,fill opacity=0.8,text opacity=1,inner sep=1.3pt,right=2pt,font=\footnotesize,text=blue!50!black]{振幅 $A=3$};
  % 峰、谷标点
  \filldraw[blue] (axis cs:2.356194,4) circle (1.7pt);
  \filldraw[blue] (axis cs:0.7853982,-2) circle (1.7pt);
  % 最大值/最小值标注
  \draw (axis cs:2.356194,4) node[above=1pt,font=\footnotesize,text=blue!50!black]{最大值 $4$};
  \draw (axis cs:0.7853982,-2) node[below=1pt,font=\footnotesize,text=blue!50!black]{最小值 $-2$};

  % ===== 竖直平移 k=1 标注（从 0 抬到中线）=====
  \draw[<->,red!70!black,thick] (axis cs:-1.0,0) -- (axis cs:-1.0,1);
  \draw (axis cs:-1.0,0.5) node[right=2pt,fill=white,fill opacity=0.82,text opacity=1,inner sep=1.2pt,font=\footnotesize,text=red!60!black]{竖直平移 $k=1$};

  % ===== 周期 T = 2pi/omega = pi 标注（沿中线，两次过中线上升点之间）=====
  % 目标曲线过中线向上的点：2x - pi = 0 => x = pi/2 = 1.5708；下一次 x = pi/2 + pi = 3pi/2... 用一个完整周期 1.5708 -> 1.5708+pi
  \draw[<->,orange!85!black,thick] (axis cs:1.5708,1) -- (axis cs:4.712389,1);
  \draw (axis cs:3.55,1.55) node[fill=white,fill opacity=0.82,text opacity=1,inner sep=1.3pt,font=\footnotesize,text=orange!60!black]{周期 $T=\dfrac{2\pi}{\omega}=\pi$};

  % ===== 相位/水平平移 标注 =====
  % 提取后 2(x - pi/2)，相对基准右移 pi/2。基准峰在 x=pi/2；目标对应峰... 用过中线上升点对比：
  % 基准 sin x 过原点上升点 x=0；目标过中线上升点 x=pi/2 -> 右移 pi/2
  \draw[->,violet,thick] (axis cs:0,0) to[bend left=22] (axis cs:1.5708,1);
  \draw (axis cs:1.95,-1.85) node[fill=white,fill opacity=0.82,text opacity=1,inner sep=1.3pt,font=\footnotesize,text=violet]{相位平移 $-\dfrac{\varphi}{\omega}=\dfrac{\pi}{2}$（右移）};
  \filldraw[gray] (axis cs:0,0) circle (1.4pt);
  \filldraw[blue] (axis cs:1.5708,1) circle (1.6pt);

\end{axis}
\end{tikzpicture}
\end{document}
